\chapter{Results}

\section{Factual simulation of Tropical Cyclone (TC) Senyar with RegCM5}
In Figure \ref{fig:factual_track}, we present the best factual simulation to represent TC Senyar from experiments summarized in Appendix \ref{appendix:exp}. The convection-permitting model (CPM) adds a significant value over the ERA5 reanalysis, which failed to represent the cyclone's trajectory and intensity as accurately as the IBTrACS observation.

\begin{figure}[h]
\centering
\includegraphics[width=\textwidth]{../../figs/track/regcm_track_best_2_new.png}
\caption{Track comparison between ERA5, observation (IBTrACS), and factual simulation}
\label{fig:factual_track}
\end{figure}

\section{Anthropogenic influences on Tropical Cyclone Senyar}

%\paragraph{Thermodynamic Baseline and Synoptic Fidelity}
We first verify the experimental design through near-surface air temperature and geopotential height responses. Time-mean near-surface temperature plots for the counterfactual simulations (Figure \ref{fig:t2m_change}) confirm the successful application of the warming and cooling signals. These simulations reveal model-specific climate sensitivities. For instance, the NorESM2-MM model exhibits stronger land-based warming/cooling anomalies than other models, implying the importance of local feedback mechanisms.

\begin{figure}[h]
\centering
\includegraphics[width=\linewidth]{../../figs/compare/tas_complete.png}
\caption{Difference of time-mean (25--28 November 2025) near-surface temperature between counterfactual and factual simulation}
\label{fig:t2m_change}
\end{figure}

This thermodynamic shift is further reflected in the 500 hPa geopotential height plots (Figure \ref{fig:geop_change}), which serve as a proxy for synoptic consistency. In the present-day simulation, the RegCM5 model faithfully reproduces the ERA5 reanalysis baseline, maintaining spatial structures around 5,840--5,870 meters. Under the +1.5$^\circ$ future scenario, widespread thermal expansion elevates the entire 500 hPa surface above the 5,870-meter threshold, while the -1.5$^\circ$ past scenario depresses it below 5,820 meters. Critically, the PGW framework succeeds in preserving the core synoptic features—specifically the cyclone vortex—while correctly adjusting the hydrostatic pressure heights and atmospheric column temperatures.

\begin{figure}[h]
\centering
\includegraphics[width=0.8\linewidth]{../../figs/dynamical/zg500_all_model.png}
\caption{Time-mean of 500 hPa geopotential height at 25 November 2025. Ensemble mean is used for counterfactuals.}
\label{fig:geop_change}
\end{figure}

%\paragraph{Track Robustness and Peak Wind Intensification}
Next, we evaluate the representation of Tropical Cyclone Senyar across all experiments. Across all climate scenarios—past, present, and future—the trajectory of TC Senyar remains robust, confirming that the storm's path is not significantly altered by background warming. While some counterfactuals fail to represent proper track after landfall in Sumatra, each still shows reasonable sea level pressure of the tropical cyclone lifetime and maximum wind speed.

While the general track and overall intensity (in terms of minimum sea level pressure) do not show significant shifts, a distinct anthropogenic influence can be identified in the wind intensity. In the warmer future climate, the results indicate that wind intensity consistently increases at peak times. This is consistent with previous study that found robust increase in wind speed of TCs \citep{patricola2018anthropogenic}. This suggests that while the 2025 event was primarily a hydrological disaster, similar equatorial cyclones in a warmer world may evolve into compounding hazards with more destructive wind profiles upon landfall.

\begin{figure}[h]
\centering
\includegraphics[width=\linewidth]{../../figs/track/regcm_track_final_0_3hr.png}
\caption{Simulated TC Senyar track and intensity. (a) For track, ensemble median of counterfactuals track is used. Tracks from each experiment are plotted as shaded lines. (b, c) For intensity, ensemble mean $\pm$ standard deviation of the counterfactuals is used.}
\label{fig:track_exp_all}
\end{figure}

% To understand the synoptic evolution of the factual simulation, daily view of sea level pressure, total precipitation, and near-surface wind speed are presented in Figure \ref{fig:synoptic_factual}.

% \begin{figure}[h]
% \centering
% \includegraphics[width=\linewidth]{../figs/synoptic_factual/synoptic.png}
% \caption{Evolution of daily precipitiation from factual simulation. Mean sea level pressure is presented as contour lines. Daily mean near-surface wind speed is presented as arrows.}
% \label{fig:synoptic_factual}.
% \end{figure}

\section{Anthropogenic influences on rainfall associated with Tropical Cyclone Senyar}

Before assessing the human fingerprint on the disaster, it is essential to verify the model's ability to represent the 2025 extreme event. A comparison of accumulated precipitation during the TC Senyar lifetime (25--28 November 2025) between the factual (present-day) simulation, the MSWEP v2 observed gridded dataset, ERA5 reanalysis, and local stations demonstrates the added value of the high-resolution RegCM5 model (Figure \ref{fig:compare_observed_pr}). While the ERA5 reanalysis captures the broad synoptic evolution of the storm, it tends to underestimate the magnitude and spatial concentration of local extremes in the Aceh province. In contrast, the factual simulation more closely aligns with the spatial distribution of extremes observed in the MSWEP dataset, successfully capturing the intense precipitation core over northern Sumatra. Although the model exhibits some discrepancies in representing the precise daily evolution of rainfall compared to local station data (See Appendix \ref{appendix:daily_evolution}), its superior performance in reproducing the spatial extent of heavy rainfall relative to ERA5 provides a foundation for further attribution and projection analyses.

\begin{figure}[h]
\centering
\includegraphics[width=0.95\linewidth]{../../figs/compare/finals/observed_pr.png}
\caption{Comparison of accumulated precipitation during TC Senyar lifetime between factual (present), MSWEP, and ERA5. Corresponding values from local stations are plotted as scatter points. Factual simulation with RegCM improves spatial distribution of extremes compared to ERA5.}
\label{fig:compare_observed_pr}
\end{figure}

Comparing the present-day simulation to the pre-industrial counterfactual reveals that human-induced climate change has already significantly amplified the 2025 event. In the present climate, the total rainfall in Aceh increased by 10\%, matched by a 6\% increase in heavy rainfall intensity (99.9th percentile of hourly intensity). The most striking anthropogenic influence, however, is the spatial expansion of the hazard: the rainfall area exceeding the 300 mm threshold grew by 64\% compared to the cooler pre-industrial world. This expansion indicates that historical warming has not only made the storm wetter but has drastically increased the geographic area exposed to extreme hydrological triggers. These historical increases generally align with the Clausius-Clapeyron (CC) relationship, where a warmer atmosphere holds approximately 7\% more moisture per degree of warming.

\begin{figure}[ht]
    \centering
    \begin{subfigure}{0.8\textwidth}
        \centering
        \includegraphics[width=\textwidth]{../../figs/compare/finals/map_pr_mean_with_impact.png}
        % \caption{}
    \end{subfigure}
    % \vspace{1em} % optional spacing
    \begin{subfigure}{0.8\textwidth}
        \centering
        \includegraphics[width=\textwidth]{../../figs/compare/finals/dist_pr_mean.png}
        % \caption{}
    \end{subfigure}
    \caption{(Top) Spatial map of accumulated precipitation between peak days. Ensemble mean is used for counterfactuals. Triangles are impact markers. (Bottom) Distribution change based on pooled hourly intensity from Aceh Province during peak days (orange line).}
    \label{fig:compare_pr}
\end{figure}

Projecting the 2025 event into a +1.5$^\circ$ warmer future introduces a more complex changes. While heavy rainfall intensity is projected to continue its increase (+17\%), this intensification does not lead to a larger total rainfall volume or a broader impact zone. Instead, the results indicate a reversal of the historical trend: total rainfall is projected to decrease by 9\%, and the rainfall area is expected to contract by 15\%. This suggests that in a warmer future, atmospheric moisture is not distributed equally across all events; rather, rainfall becomes increasingly concentrated into fewer but more violent bursts. While the overall footprint of the storm may shrink, the localized core of the disaster becomes more dangerous as extreme intensity begins to exceed Clausius-Clapeyron scaling, even as lighter rainfall events decrease in frequency.

The evolution of area-averaged hourly rainfall and its corresponding cumulative totals in Aceh Province provide visual confirmation of the discussed changes. The cumulative hyetograph (Figure \ref{fig:hyetograph}) shows that with present-day totals ending higher than those in the pre-industrial (-1.5$^\circ$) scenario. In contrast, the future (+1.5$^\circ$) cumulative trajectory eventually falls below present-day levels, supporting the projected 9\% decrease in total rainfall. %Physically, this indicates that while a warmer future atmosphere holds more moisture, the lack of consistently favorable lifting mechanisms—such as atmospheric circulation or storm dynamics—results in a less efficient conversion of that moisture into surface precipitation over the storm's duration

\begin{figure}[h]
\centering
\includegraphics[width=0.9\linewidth]{../../figs/compare/finals/hyetograph_event.png}
\caption{(a) Evolution of area-averaged hourly rainfall in Aceh Province with its ensemble mean $\pm$ standard deviation for counterfactuals, while (b) shows the corresponding cumulative rainfall.}
\label{fig:hyetograph}
\end{figure}

The physical basis for these changes lies in the efficiency with which the atmosphere converts moisture into rain. Although a warmer future atmosphere holds more moisture, every precipitation event does not necessarily produce more rainfall. The dynamical factors that control the efficiency of how moisture is converted to precipitation will be discussed in the next subchapter. In the case of future projection of Aceh, the intensification of extreme intensity suggests that the localized convective process is more efficient. This changing feature implies that future adaptation strategies must prioritize resilience against concentrated, high-magnitude bursts that can trigger devastating flash floods and landslides, even within a smaller total disaster footprint.

\section{Anthropogenic influences on the physical drivers of extreme rainfall}

In understanding changes in extreme precipitation, it is common to decompose the contributions from thermodynamics and dynamics \citep{pfahl2017understanding} while both interact in a feedback especially in the tropics \citep{neelin2022precipitation}. Thus, to explain the physical mechanisms behind the 2025 Sumatra disaster, this subchapter analyzes the anthropogenic influences in the thermodynamic and dynamical drivers that control the extreme rainfall. By quantifying changes in moisture supply, atmospheric instability, and convective strength, we can identify how human-caused warming has intensified the extreme rainfall associated with TC Senyar and how it changes in the future.

\subsection{Changes in moisture supply: Evaporation and IVT}
The intensification of the 2025 extreme rainfall begins with an increased supply of moisture from the surrounding oceans. Analysis of evaporation flux in Figure \ref{fig:compare_evaporation} shows a consistent median increase of +4\% in both the historical attribution (past-to-present) and future projection scenarios (present-to-future). This steady increase in surface moisture is complemented by significant changes in low-level integrated vapor transport (IVT) over the Malacca Strait. The IVT spatial maps and exceedance probability distributions (Figure \ref{fig:compare_ivt}) reveal a critical shift in moisture dynamics.

In a warmer climate, the extreme tail of IVT (values exceeding 800 kg m$^{-1}$ s$^{-1}$) increases substantially at a rate of 5\% to 12\% per degree of warming, closely aligning with Clausius-Clapeyron scaling. This intensification is concentrated specifically within the storm’s core moisture corridor over the Malacca Strait, while surrounding areas experience a decrease in weak moisture transport. This suggests that anthropogenic warming converges more extreme moisture directly into the cyclone’s circulation.

\begin{figure}[ht]
    \centering
    \begin{subfigure}{0.8\textwidth}
        \centering
        \includegraphics[width=\textwidth]{../../figs/compare/finals/map_evs.png}
        % \caption{}
    \end{subfigure}
    % \vspace{1em} % optional spacing
    \begin{subfigure}{0.8\textwidth}
        \centering
        \includegraphics[width=\textwidth]{../../figs/compare/finals/dist_evs.png}
        % \caption{}
    \end{subfigure}
    \caption{(Top) Spatial map of time-mean evaporation flux during peak days. Ensemble mean is used for counterfactuals. (Bottom) Distribution change based on pooled hourly values from whole shown domain.}
    \label{fig:compare_evaporation}
\end{figure}

\begin{figure}[ht]
    \centering
    \begin{subfigure}{0.8\textwidth}
        \centering
        \includegraphics[width=\textwidth]{../../figs/compare/finals/map_ivt.png}
        % \caption{}
    \end{subfigure}
    % \vspace{1em} % optional spacing
    \begin{subfigure}{0.8\textwidth}
        \centering
        \includegraphics[width=\textwidth]{../../figs/compare/finals/dist_ivt.png}
        % \caption{}
    \end{subfigure}
    \caption{(Top) Spatial map of time-mean low-level IVT during peak days. Low-level IVT is obtained by integration between surface and 700 hPa. Ensemble mean is used for counterfactuals. (Bottom) Distribution change based on pooled six-hourly values from whole shown domain.}
    \label{fig:compare_ivt}
\end{figure}

\subsection{Changes in atmospheric instability and potential energy}
The increased moisture supply is coupled with a dramatic rise in convective available potential energy (CAPE), the primary measure of atmospheric instability, seen in Figure \ref{fig:compare_cape}. Across both climate scenarios, median CAPE values increase by +23\%--28\%, with high-CAPE regions (>1000 J/kg) expanding significantly in both spatial area and intensity along the storm corridor (the Malacca Strait). These findings support previous research \citep{seeley2015does}, showing that CAPE increases as the climate warms, providing the potential required for severe convective storms. This surge in potential energy ensures that when moisture is lifted within the cyclone's core, the resulting convection is significantly more violent than it would have been in a cooler climate.

\begin{figure}[ht]
    \centering
    \begin{subfigure}{0.8\textwidth}
        \centering
        \includegraphics[width=\textwidth]{../../figs/compare/finals/map_cape.png}
        % \caption{}
    \end{subfigure}
    % \vspace{1em} % optional spacing
    \begin{subfigure}{0.8\textwidth}
        \centering
        \includegraphics[width=\textwidth]{../../figs/compare/finals/dist_cape.png}
        % \caption{}
    \end{subfigure}
    \caption{(Top) Spatial map of time-mean CAPE during peak days. Ensemble mean is used for counterfactuals. (Bottom) Distribution change based on pooled hourly values from whole shown domain.}
    \label{fig:compare_cape}
\end{figure}

\subsection{Changes in physical mechanisms: Latent Heating and Updrafts}
The final link in the convective process is the conversion of potential energy into vertical motion and intense rainfall through dynamical feedbacks. As moisture condenses within the storm's core, it releases more latent heat, which further amplifies the peak diabatic heating rates, up to 17\% per degree of warming in future scenarios (Figure \ref{fig:compare_diab}). This localized release of energy increases in-cloud buoyancy, which directly accelerates vertical motion (Figure \ref{fig:compare_wa}).

Consequently, the 99.9th percentile of maximum updraft velocity shows an increase of +12\% in the present climate and +15\% future. This dynamical intensification explains the complex intensity-area change: even as the broader regional rainfall area contracts, the intensified convective core becomes more efficient at converting moisture to concentrated, high-magnitude bursts.

\begin{figure}[ht]
    \centering
    \begin{subfigure}{0.8\textwidth}
        \centering
        \includegraphics[width=\textwidth]{../../figs/compare/finals/map_diab.png}
        % \caption{}
    \end{subfigure}
    % \vspace{1em} % optional spacing
    \begin{subfigure}{0.8\textwidth}
        \centering
        \includegraphics[width=\textwidth]{../../figs/compare/finals/dist_diab.png}
        % \caption{}
    \end{subfigure}

    \caption{(Top) Spatial map of time-maximum of maximum temperature tendency due to latent heat exchange during peak days. Maximum temperature tendency is obtained by taking vertical maxima over each time step. Ensemble mean is used for counterfactuals. (Bottom) Distribution change is based on pooled six-hourly values from whole shown domain.}
    \label{fig:compare_diab}
\end{figure}

\begin{figure}[ht]
    \centering
    \begin{subfigure}{0.8\textwidth}
        \centering
        \includegraphics[width=\textwidth]{../../figs/compare/finals/map_wa_max.png}
        % \caption{}
    \end{subfigure}
    % \vspace{1em} % optional spacing
    \begin{subfigure}{0.8\textwidth}
        \centering
        \includegraphics[width=\textwidth]{../../figs/compare/finals/dist_wa.png}
        % \caption{}
    \end{subfigure}

    \caption{(Top) Spatial map of time-maximum of maximum vertical velocity during peak days. Maximum vertical velocity is obtained by taking vertical maxima over each time step. Ensemble maxima is used for counterfactuals. (Bottom) Distribution change is based on pooled six-hourly values from whole shown domain.}
    \label{fig:compare_wa}
\end{figure}